\section{结论}
\label{sec:conclusion}

本文结合相匹配初态、反射对称性约束的变分拟设与关联量测量，研究了二聚化 XXZ 链的低能态制备。归一化部分反射标记给出 $L=8$ 的有限尺寸相图，区分平庸、拓扑与 AFM 区域；$\delta=-2$ 线上的对称性扇区能隙分析显示，随链长增加，$s=0.5$ 附近的能隙减小。

在百花处理器上，候选子图通过簇态稳定子与读出保真度进行筛选，测量分布采用逐批标定的张量积读出模型修正。真机与模拟机的 $Q$ 曲线在 $\delta=-2$ 扫描线上呈现相近的参数依赖；在 $\delta=0.85$ 的拓扑区域，预期方向上没有可分辨的正 string 信号：$s=0.99$ 处 pooled $\hat O_{\rm str}$ 为 $-0.044 \pm 0.005$，与模拟值 $+0.95$ 符号相反，其余多个点同样显著偏离零；真机 $Q$ 曲线在 $s\to1$ 时与零在窗口不确定度内一致。

这些结果展示了在同一框架下结合相结构、对称性约束和直接关联测量的实验流程，并将拓扑区域的 string 信号保留确定为后续改进的重要目标。进一步研究可扩展链长与参数扫描，并结合关联读出模型考察噪声对非局域关联的影响。
